\chapter{Mirror TBA and the Y-System}
\label{ch:planar-n4-mirror-tba-y-system}

\ConventionsLorentzian

The asymptotic Bethe ansatz determines long-chain states before wrapping
effects.  Finite single traces require a finite-size formalism.  The planar
\(\mathcal N=4\) construction uses a mirror model obtained by exchanging
worldsheet space and Euclidean time in the integrable spin-chain scattering
problem.  This chapter states the logical structure without pretending that
the mirror thermodynamic Bethe ansatz is a theorem of four-dimensional
constructive QFT.

\section{Mirror Transformation}

For a physical magnon with energy \(E(p)\) and momentum \(p\), the mirror
variables are defined by analytic continuation
\[
  \widetilde E=\ii p,\qquad
  \widetilde p=\ii E .
\]
Since
\[
  E(p)=\sqrt{1+16g^2\sin^2\frac p2},
\]
the mirror dispersion is not the same function of \(\widetilde p\) as the
physical dispersion.  This is the first major difference from relativistic
two-dimensional QFT, where the mirror transformation is a rapidity shift.

Bound states of \(Q\) physical magnons are described by pairs
\[
  x_Q^+,\qquad x_Q^-,
\]
obeying the bound-state shortening relation
\[
  x_Q^++\frac1{x_Q^+}-x_Q^- -\frac1{x_Q^-}
  =
  \frac{\ii Q}{g}.
\]
In the mirror region the branch is chosen so that the mirror energy
\(\widetilde E_Q(u)\) is positive for real mirror rapidity \(u\).

\section{Thermodynamic Bethe Ansatz}

Let \(A,B\) range over the mirror Bethe-string species: physical
\(Q\)-particle bound states, auxiliary \(y\)-roots, and auxiliary \(w\)-string
families for the two \(\mathfrak{su}(2|2)\) wings.  Denote the mirror
pseudoenergy by \(\epsilon_A(u)\) and set
\[
  Y_A(u)=\ee^{-\epsilon_A(u)}.
\]
The TBA equations have the form
\begin{equation}
  \epsilon_A(u)
  =
  L\,\widetilde E_A(u)-\mu_A
  -
  \sum_B\int_{\mathcal C_B}\dd v\,
  K_{BA}(v,u)\log\!\left(1+Y_B(v)\right),
  \label{eq:planar-n4-mirror-tba-general}
\end{equation}
with kernels
\[
  K_{BA}(v,u)
  =
  \frac{1}{2\pi\ii}\frac{\partial}{\partial v}
  \log S_{BA}^{\rm mirror}(v,u).
\]
The contours \(\mathcal C_B\), signs for fermionic auxiliary roots, and
chemical potentials \(\mu_A\) are part of the mirror-model definition.  They
are not decorations: changing them changes the finite-size answer.

For a state with physical Bethe roots \(u_1,\ldots,u_K\), the excited-state
TBA is obtained by contour deformation.  The physical roots satisfy exact
Bethe equations of the form
\[
  1+Y_{\bullet_1}^{\rm phys}(u_j)=0,
  \qquad j=1,\ldots,K,
\]
and the energy is
\begin{equation}
  \Delta-J
  =
  \sum_{j=1}^K E(u_j)
  -
  \sum_{Q=1}^\infty
  \int_{\R}\frac{\dd u}{2\pi}\,
  \frac{\dd\widetilde p_Q}{\dd u}
  \log(1+Y_{\bullet_Q}(u)).
  \label{eq:planar-n4-excited-tba-energy}
\end{equation}
The first term is the asymptotic magnon energy.  The integral term is the
finite-size wrapping correction from the thermal mirror bath.

\section{Y-System}

The TBA equations imply functional relations among Y-functions after applying
kernel identities and shifting the spectral parameter.  On the interior of
the planar \(\mathcal N=4\) T-hook, the relation has the form
\begin{equation}
  Y_{a,s}^+(u)Y_{a,s}^-(u)
  =
  \frac{(1+Y_{a,s+1}(u))(1+Y_{a,s-1}(u))}
       {(1+Y_{a+1,s}(u)^{-1})(1+Y_{a-1,s}(u)^{-1})},
  \label{eq:planar-n4-y-system}
\end{equation}
where
\[
  f^\pm(u)=f(u\pm\ii/2).
\]
The indices \((a,s)\) lie in the T-hook domain determined by the two
\(\mathfrak{su}(2|2)\) wings.  Boundary nodes require modified relations.

\begin{warning}[Y-system data are not only equations]
\label{warn:y-system-not-complete}
The functional relations \eqref{eq:planar-n4-y-system} do not by themselves
determine a spectrum.  A solution also requires analyticity domains,
large-\(u\) asymptotics, branch-cut discontinuity relations, exact Bethe-root
conditions, and state labels.  Treating the Y-system equations alone as a
definition of the theory discards the information that distinguishes physical
states.
\end{warning}

\section{Konishi as the First Wrapping Test}

The Konishi multiplet contains an \(SL(2)\) descendant represented by a
length-two state with two derivative magnons.  In the asymptotic Bethe ansatz,
the two physical roots are opposite,
\[
  u_1=-u_2.
\]
At weak coupling the one-loop solution is
\[
  u_1=\frac{1}{2\sqrt3},\qquad
  u_2=-\frac{1}{2\sqrt3},
\]
in the normalization of the rational Bethe equations.  The asymptotic
dimension matches direct perturbative gauge theory through the orders before
wrapping appears.  The first wrapping correction is supplied by the mirror
integral term in \eqref{eq:planar-n4-excited-tba-energy}.  Conceptually, this
is the finite-length correction to the trace, not a correction to the
two-magnon dispersion relation.

The exact finite-coupling Konishi dimension is obtained by solving the
excited-state TBA or, more efficiently, the quantum spectral curve of the next
chapter.  In this monograph the TBA formula is used as a logically explicit
construction of the finite-size correction: every input is an S-matrix kernel,
a contour, a chemical potential, or a Bethe-root condition.

\section{Relation to Relativistic Integrable QFT}

The nested TBA of
Chapter~\ref{ch:integrable-nested-tba-tq-separation-variables} supplies the
finite-density variational logic.  The present system differs in three ways.
First, the dispersion is not relativistic.  Second, the mirror theory has an
infinite family of bound-state and auxiliary nodes.  Third, the spectral
problem is a gauge-theory local-operator problem, so the circle length is the
trace length \(L\) or a charge shifted from it by sector-dependent
conventions.  These differences are why the planar \(\mathcal N=4\) TBA is
placed in the supersymmetric gauge-theory volume rather than in the
relativistic integrable-QFT volume.
